\section*{Abstract}
\addcontentsline{toc}{section}{Abstract}

Frontier language models are now the default backbone for agentic systems, but
their dependence on external APIs limits their use in offline,
privacy-sensitive, and high-volume deployments. Small language models can run on
consumer hardware and avoid these constraints, but they often fail at the
structured output and tool-calling tasks that agentic pipelines require. This
thesis examines how far optimization techniques can move a small model toward
reliable tool calling, and which accuracy cost or benefit each technique
introduces.

The study uses an incremental ablation over the Qwen~2.5 model family
(0.5B--7B), with \texttt{Qwen2.5-7B-Instruct} as the primary model. The models
are evaluated on four BFCL v4 categories (\texttt{simple\_python},
\texttt{multiple}, \texttt{parallel}, and \texttt{parallel\_multiple}) and on
$\tau$-bench retail. A minimal model-agnostic baseline that does not use the
% TODO(#redteam): framing. The 1.5\% / "77--96 pp below frontier" gap is the
% strict whole-completion-parse floor. The same 7B run under the lenient parser
% scores 62.00\%; report the strict floor as a floor, and consider stating the
% same-parser figure so the headline gap is not read as an intrinsic deficit.
model's native tool-calling template scores 1.5\% on
\texttt{simple\_python}, which is 77--96 percentage points below frontier
models on the same category (78.75\%--97.75\%). A control configuration using
the model's native template reaches 96.00\%, inside the frontier band. The
unoptimized gap therefore measures the cost of a generic integration rather
than the model's capability ceiling.

Constrained decoding is the dominant intervention. It raises accuracy to
72.75\% without modifying model weights or generating additional tokens, while
reducing mean per-call latency from 6.995~s under unconstrained generation to
0.878~s. AWQ INT4 quantization preserves most of this accuracy at 72.25\%
while reducing memory from 14.25~GiB to 5.20~GiB. In contrast, few-shot
prompting and chain-of-thought reasoning reduce accuracy by 2.5 and 6.75
percentage points relative to the constrained-decoding configuration each
extends. Retrieval-augmented generation reduces accuracy by 24.5 percentage
points, but in an open-catalogue setting that replaces the single oracle
candidate with five retrieved candidates; the drop reflects candidate
disambiguation (retrieval recall@5 is 97.2\%) rather than retrieval overhead on
a fixed task. Format-aligned LoRA
fine-tuning reaches 76.75\%, and 74.25\% after subsequent quantization at the
same memory footprint.

The strongest no-training result is schema-enriched constrained decoding
(CD+schema), which adds parameter descriptions, enumeration values, and default
values to the argument extraction prompt while keeping the FSM structural
constraint unchanged. CD+schema reaches 89.0\% ($+$16.25~pp over plain CD,
McNemar $p < 0.00001$), placing the 7B model inside the frontier range on the
single-call, single-candidate BFCL category without modifying model weights. The
unoptimized failure mode is therefore primarily a format-compliance gap: 394 of
400 baseline failures are malformed outputs rather than wrong answers.
Constrained decoding removes that gap structurally, and schema enrichment
recovers much of the remaining loss by restoring schema information discarded by
the generic prompt. The result does not imply end-to-end agent reliability.
Accuracy falls to 4.3\% on $\tau$-bench retail, and BFCL \texttt{parallel}
scores 0\% under the initial single-call schema, though extending the schema to
joint multi-call generation lifts it to 79\% at 7B. The persistent
$\tau$-bench gap shows that multi-turn planning, unlike multi-call emission,
requires more than the optimization techniques studied here.
